\chapter{Практическая часть}

\section*{Программа <<Телефонный справочник>>}

Разработать свою программу --- <<Телефонный справочник>>. 
Владельцы могут иметь несколько телефонов и автомобилей (Факты).
Протестировать работу программы, используя разные вопросы. 
\begin{itemize}
	\item <<Телефонный справочник>>: Фамилия, №тел, Адрес – структура (Город, Улица, №дома, №кв);
	\item <<Автомобиль>>: Фамилия, Марка, Цвет, Стоимость, Номер.
\end{itemize}

\begin{lstlisting}[caption=Программа <<Телефонный справочник>> (Начало)]
domains
	name, surname=string.
	phone_list=string*.
	city, street=string.
	house, flat=integer.
	address=address_data(city, street, house, flat).
	price, car_number=integer.
	colour, brand=string.

predicates
	phone_record(surname, name, address, phone_list).
	car(surname, name, brand, colour, price, car_number).
	question(surname, city, phone_list, brand, colour).
	
clauses
	car("Иванов", "Александр", "Mercedes", "black", 200000, 2313).
	car("Иванов", "Виктор", "Mercedes", "yellow", 200000, 4312).
	car("Сидоров", "Александр", "Lada", "red", 50000, 5555).
	car("Жаров", "Андрей", "Toyta", "green", 100000, 9999).
	car("Петров", "Евгений", "Audi", "gray", 600000, 8888).
	car("Петров", "Евгений", "Mercedes", "black", 500000, 7788).
\end{lstlisting}

\clearpage
	
\begin{lstlisting}[caption=Программа <<Телефонный справочник>> (Продолжение)]
	phone_record("Иванов", "Александр", address_data("Москва", "Абельмановская улица", 10, 20), ["+89280202902", "+891904802", "+891904805"]).
	phone_record("Иванов", "Виктор", address_data("Новосибирск", "Авиаматорная улица", 1, 10), ["8456372"]).
	phone_record("Сидоров", "Александр", address_data("Москва", "Анненский проезд", 3, 15), ["8994527"]).
	phone_record("Жаров", "Андрей", address_data("Москва", "улица Носова", 1, 16), ["8994558"]).
	phone_record("Петров", "Евгений", address_data("Москва", "улица Петра", 5, 10), ["6994566"]).
	
	question(Surname, City, Phone, Brand, Colour) :-
	phone_record(Surname, _, address_data(City, _, _, _), Phone),
	car(Surname, _, Brand, Colour, _, _).
\end{lstlisting}

Задать вопрос такой чтобы по Марке и Цвету автомобиля найти Фамилию, Город, Телефон.
\begin{lstlisting}[caption=Программа <<Телефонный справочник>> (Вопрос)]
goal
	% Какую фамилию, номер, город, имеет владелец автомобиля конкретной марки и цвета автомобиля?
	question(Surname, City, Phone, "Mercedes", "black").
\end{lstlisting}

Другие вопросы:
\begin{lstlisting}
	question("Иванов", "Новосибирск", ["8456372"], Brand, Colour).
	
	phone_record(Surname, _, address_data("Москва", _, _, _), _), car(Surname, _, Mark, _, _, CarNumber).
	
	car("Жаров", _, "Toyta", _, _, _).
	
	car(Surname, Name,  "Mercedes", _, _, _).
\end{lstlisting}